%% bare_jrnl_transmag.tex
%% V1.4b
%% 2015/08/26
%% by Michael Shell
%% see http://www.michaelshell.org/
%% for current contact information.
%%
%% This is a skeleton file demonstrating the use of IEEEtran.cls
%% (requires IEEEtran.cls version 1.8b or later) with an IEEE 
%% Transactions on Magnetics journal paper.
%%
%% Support sites:
%% http://www.michaelshell.org/tex/ieeetran/
%% http://www.ctan.org/pkg/ieeetran
%% and
%% http://www.ieee.org/

%%*************************************************************************
%% Legal Notice:
%% This code is offered as-is without any warranty either expressed or
%% implied; without even the implied warranty of MERCHANTABILITY or
%% FITNESS FOR A PARTICULAR PURPOSE! 
%% User assumes all risk.
%% In no event shall the IEEE or any contributor to this code be liable for
%% any damages or losses, including, but not limited to, incidental,
%% consequential, or any other damages, resulting from the use or misuse
%% of any information contained here.
%%
%% All comments are the opinions of their respective authors and are not
%% necessarily endorsed by the IEEE.
%%
%% This work is distributed under the LaTeX Project Public License (LPPL)
%% ( http://www.latex-project.org/ ) version 1.3, and may be freely used,
%% distributed and modified. A copy of the LPPL, version 1.3, is included
%% in the base LaTeX documentation of all distributions of LaTeX released
%% 2003/12/01 or later.
%% Retain all contribution notices and credits.
%% ** Modified files should be clearly indicated as such, including  **
%% ** renaming them and changing author support contact information. **
%%*************************************************************************


% *** Authors should verify (and, if needed, correct) their LaTeX system  ***
% *** with the testflow diagnostic prior to trusting their LaTeX platform ***
% *** with production work. The IEEE's font choices and paper sizes can   ***
% *** trigger bugs that do not appear when using other class files.       ***                          ***
% The testflow support page is at:
% http://www.michaelshell.org/tex/testflow/



\documentclass[journal,transmag]{IEEEtran}
%
% If IEEEtran.cls has not been installed into the LaTeX system files,
% manually specify the path to it like:
% \documentclass[journal]{../sty/IEEEtran}





% Some very useful LaTeX packages include:
% (uncomment the ones you want to load)


% *** MISC UTILITY PACKAGES ***
%
%\usepackage{ifpdf}
% Heiko Oberdiek's ifpdf.sty is very useful if you need conditional
% compilation based on whether the output is pdf or dvi.
% usage:
% \ifpdf
%   % pdf code
% \else
%   % dvi code
% \fi
% The latest version of ifpdf.sty can be obtained from:
% http://www.ctan.org/pkg/ifpdf
% Also, note that IEEEtran.cls V1.7 and later provides a builtin
% \ifCLASSINFOpdf conditional that works the same way.
% When switching from latex to pdflatex and vice-versa, the compiler may
% have to be run twice to clear warning/error messages.






% *** CITATION PACKAGES ***
%
%\usepackage{cite}
% cite.sty was written by Donald Arseneau
% V1.6 and later of IEEEtran pre-defines the format of the cite.sty package
% \cite{} output to follow that of the IEEE. Loading the cite package will
% result in citation numbers being automatically sorted and properly
% "compressed/ranged". e.g., [1], [9], [2], [7], [5], [6] without using
% cite.sty will become [1], [2], [5]--[7], [9] using cite.sty. cite.sty's
% \cite will automatically add leading space, if needed. Use cite.sty's
% noadjust option (cite.sty V3.8 and later) if you want to turn this off
% such as if a citation ever needs to be enclosed in parenthesis.
% cite.sty is already installed on most LaTeX systems. Be sure and use
% version 5.0 (2009-03-20) and later if using hyperref.sty.
% The latest version can be obtained at:
% http://www.ctan.org/pkg/cite
% The documentation is contained in the cite.sty file itself.






% *** GRAPHICS RELATED PACKAGES ***
%
\ifCLASSINFOpdf
  % \usepackage[pdftex]{graphicx}
  % declare the path(s) where your graphic files are
  % \graphicspath{{../pdf/}{../jpeg/}}
  % and their extensions so you won't have to specify these with
  % every instance of \includegraphics
  % \DeclareGraphicsExtensions{.pdf,.jpeg,.png}
\else
  % or other class option (dvipsone, dvipdf, if not using dvips). graphicx
  % will default to the driver specified in the system graphics.cfg if no
  % driver is specified.
  % \usepackage[dvips]{graphicx}
  % declare the path(s) where your graphic files are
  % \graphicspath{{../eps/}}
  % and their extensions so you won't have to specify these with
  % every instance of \includegraphics
  % \DeclareGraphicsExtensions{.eps}
\fi
% graphicx was written by David Carlisle and Sebastian Rahtz. It is
% required if you want graphics, photos, etc. graphicx.sty is already
% installed on most LaTeX systems. The latest version and documentation
% can be obtained at: 
% http://www.ctan.org/pkg/graphicx
% Another good source of documentation is "Using Imported Graphics in
% LaTeX2e" by Keith Reckdahl which can be found at:
% http://www.ctan.org/pkg/epslatex
%
% latex, and pdflatex in dvi mode, support graphics in encapsulated
% postscript (.eps) format. pdflatex in pdf mode supports graphics
% in .pdf, .jpeg, .png and .mps (metapost) formats. Users should ensure
% that all non-photo figures use a vector format (.eps, .pdf, .mps) and
% not a bitmapped formats (.jpeg, .png). The IEEE frowns on bitmapped formats
% which can result in "jaggedy"/blurry rendering of lines and letters as
% well as large increases in file sizes.
%
% You can find documentation about the pdfTeX application at:
% http://www.tug.org/applications/pdftex




% *** MATH PACKAGES ***
%
\usepackage{amsmath}
% A popular package from the American Mathematical Society that provides
% many useful and powerful commands for dealing with mathematics.
%
% Note that the amsmath package sets \interdisplaylinepenalty to 10000
% thus preventing page breaks from occurring within multiline equations. Use:
%\interdisplaylinepenalty=2500
% after loading amsmath to restore such page breaks as IEEEtran.cls normally
% does. amsmath.sty is already installed on most LaTeX systems. The latest
% version and documentation can be obtained at:
% http://www.ctan.org/pkg/amsmath





% *** SPECIALIZED LIST PACKAGES ***
%
%\usepackage{algorithmic}
% algorithmic.sty was written by Peter Williams and Rogerio Brito.
% This package provides an algorithmic environment fo describing algorithms.
% You can use the algorithmic environment in-text or within a figure
% environment to provide for a floating algorithm. Do NOT use the algorithm
% floating environment provided by algorithm.sty (by the same authors) or
% algorithm2e.sty (by Christophe Fiorio) as the IEEE does not use dedicated
% algorithm float types and packages that provide these will not provide
% correct IEEE style captions. The latest version and documentation of
% algorithmic.sty can be obtained at:
% http://www.ctan.org/pkg/algorithms
% Also of interest may be the (relatively newer and more customizable)
% algorithmicx.sty package by Szasz Janos:
% http://www.ctan.org/pkg/algorithmicx




% *** ALIGNMENT PACKAGES ***
%
%\usepackage{array}
% Frank Mittelbach's and David Carlisle's array.sty patches and improves
% the standard LaTeX2e array and tabular environments to provide better
% appearance and additional user controls. As the default LaTeX2e table
% generation code is lacking to the point of almost being broken with
% respect to the quality of the end results, all users are strongly
% advised to use an enhanced (at the very least that provided by array.sty)
% set of table tools. array.sty is already installed on most systems. The
% latest version and documentation can be obtained at:
% http://www.ctan.org/pkg/array


% IEEEtran contains the IEEEeqnarray family of commands that can be used to
% generate multiline equations as well as matrices, tables, etc., of high
% quality.




% *** SUBFIGURE PACKAGES ***
%\ifCLASSOPTIONcompsoc
%  \usepackage[caption=false,font=normalsize,labelfont=sf,textfont=sf]{subfig}
%\else
%  \usepackage[caption=false,font=footnotesize]{subfig}
%\fi
% subfig.sty, written by Steven Douglas Cochran, is the modern replacement
% for subfigure.sty, the latter of which is no longer maintained and is
% incompatible with some LaTeX packages including fixltx2e. However,
% subfig.sty requires and automatically loads Axel Sommerfeldt's caption.sty
% which will override IEEEtran.cls' handling of captions and this will result
% in non-IEEE style figure/table captions. To prevent this problem, be sure
% and invoke subfig.sty's "caption=false" package option (available since
% subfig.sty version 1.3, 2005/06/28) as this is will preserve IEEEtran.cls
% handling of captions.
% Note that the Computer Society format requires a larger sans serif font
% than the serif footnote size font used in traditional IEEE formatting
% and thus the need to invoke different subfig.sty package options depending
% on whether compsoc mode has been enabled.
%
% The latest version and documentation of subfig.sty can be obtained at:
% http://www.ctan.org/pkg/subfig

\usepackage{enumitem}  
\usepackage{graphicx}


% *** FLOAT PACKAGES ***
%
%\usepackage{fixltx2e}
% fixltx2e, the successor to the earlier fix2col.sty, was written by
% Frank Mittelbach and David Carlisle. This package corrects a few problems
% in the LaTeX2e kernel, the most notable of which is that in current
% LaTeX2e releases, the ordering of single and double column floats is not
% guaranteed to be preserved. Thus, an unpatched LaTeX2e can allow a
% single column figure to be placed prior to an earlier double column
% figure.
% Be aware that LaTeX2e kernels dated 2015 and later have fixltx2e.sty's
% corrections already built into the system in which case a warning will
% be issued if an attempt is made to load fixltx2e.sty as it is no longer
% needed.
% The latest version and documentation can be found at:
% http://www.ctan.org/pkg/fixltx2e


%\usepackage{stfloats}
% stfloats.sty was written by Sigitas Tolusis. This package gives LaTeX2e
% the ability to do double column floats at the bottom of the page as well
% as the top. (e.g., "\begin{figure*}[!b]" is not normally possible in
% LaTeX2e). It also provides a command:
%\fnbelowfloat
% to enable the placement of footnotes below bottom floats (the standard
% LaTeX2e kernel puts them above bottom floats). This is an invasive package
% which rewrites many portions of the LaTeX2e float routines. It may not work
% with other packages that modify the LaTeX2e float routines. The latest
% version and documentation can be obtained at:
% http://www.ctan.org/pkg/stfloats
% Do not use the stfloats baselinefloat ability as the IEEE does not allow
% \baselineskip to stretch. Authors submitting work to the IEEE should note
% that the IEEE rarely uses double column equations and that authors should try
% to avoid such use. Do not be tempted to use the cuted.sty or midfloat.sty
% packages (also by Sigitas Tolusis) as the IEEE does not format its papers in
% such ways.
% Do not attempt to use stfloats with fixltx2e as they are incompatible.
% Instead, use Morten Hogholm'a dblfloatfix which combines the features
% of both fixltx2e and stfloats:
%
% \usepackage{dblfloatfix}
% The latest version can be found at:
% http://www.ctan.org/pkg/dblfloatfix




%\ifCLASSOPTIONcaptionsoff
%  \usepackage[nomarkers]{endfloat}
% \let\MYoriglatexcaption\caption
% \renewcommand{\caption}[2][\relax]{\MYoriglatexcaption[#2]{#2}}
%\fi
% endfloat.sty was written by James Darrell McCauley, Jeff Goldberg and 
% Axel Sommerfeldt. This package may be useful when used in conjunction with 
% IEEEtran.cls'  captionsoff option. Some IEEE journals/societies require that
% submissions have lists of figures/tables at the end of the paper and that
% figures/tables without any captions are placed on a page by themselves at
% the end of the document. If needed, the draftcls IEEEtran class option or
% \CLASSINPUTbaselinestretch interface can be used to increase the line
% spacing as well. Be sure and use the nomarkers option of endfloat to
% prevent endfloat from "marking" where the figures would have been placed
% in the text. The two hack lines of code above are a slight modification of
% that suggested by in the endfloat docs (section 8.4.1) to ensure that
% the full captions always appear in the list of figures/tables - even if
% the user used the short optional argument of \caption[]{}.
% IEEE papers do not typically make use of \caption[]'s optional argument,
% so this should not be an issue. A similar trick can be used to disable
% captions of packages such as subfig.sty that lack options to turn off
% the subcaptions:
% For subfig.sty:
% \let\MYorigsubfloat\subfloat
% \renewcommand{\subfloat}[2][\relax]{\MYorigsubfloat[]{#2}}
% However, the above trick will not work if both optional arguments of
% the \subfloat command are used. Furthermore, there needs to be a
% description of each subfigure *somewhere* and endfloat does not add
% subfigure captions to its list of figures. Thus, the best approach is to
% avoid the use of subfigure captions (many IEEE journals avoid them anyway)
% and instead reference/explain all the subfigures within the main caption.
% The latest version of endfloat.sty and its documentation can obtained at:
% http://www.ctan.org/pkg/endfloat
%
% The IEEEtran \ifCLASSOPTIONcaptionsoff conditional can also be used
% later in the document, say, to conditionally put the References on a 
% page by themselves.




% *** PDF, URL AND HYPERLINK PACKAGES ***
%
%\usepackage{url}
% url.sty was written by Donald Arseneau. It provides better support for
% handling and breaking URLs. url.sty is already installed on most LaTeX
% systems. The latest version and documentation can be obtained at:
% http://www.ctan.org/pkg/url
% Basically, \url{my_url_here}.




% *** Do not adjust lengths that control margins, column widths, etc. ***
% *** Do not use packages that alter fonts (such as pslatex).         ***
% There should be no need to do such things with IEEEtran.cls V1.6 and later.
% (Unless specifically asked to do so by the journal or conference you plan
% to submit to, of course. )


% correct bad hyphenation here
\hyphenation{op-tical net-works semi-conduc-tor}


\begin{document}

\title{2-Player UNO Game}

\author{\IEEEauthorblockN{Wesley Gunawan, Chen Yi-Shan, Alex Chang, Samyak Kakatur}}

\markboth{Final Report of ECE209AS, December~2025}%
{Gunawan \MakeLowercase{\textit{et al.}}: 2-Player UNO Game}

\IEEEtitleabstractindextext{%
\begin{abstract}
\textbf{\textit{Abstract}}: We model two-player UNO as a partially observable Markov decision process (POMDP) and study whether a unique Bayes-optimal move typically exists from the perspective of Player~1. For small reduced UNO decks, we enumerate all latent world states consistent with a given observation and compute exact action values for all legal plays. We show that scenarios with a strictly dominant action are surprisingly rare, even under perfect-information evaluation. Most belief states admit multiple actions with nearly identical expected value, implying that UNO optimality is structurally ambiguous. We further test a Monte Carlo belief-based agent and find that it reliably selects the Bayes-optimal move in those rare states where a unique optimum exists. These results highlight an inherent limitation of deriving deterministic “best play'' strategies in hidden-information games such as UNO.
\end{abstract}

\begin{IEEEkeywords}
UNO, POMDP, partially observable games, belief state, Monte Carlo planning.
\end{IEEEkeywords}}

\maketitle
\IEEEdisplaynontitleabstractindextext

\section{Introduction}

\IEEEPARstart{U}{NO} is a partially observable, stochastic card game in which determining optimal play is challenging due to hidden information and uncertain outcomes. In this project we investigate a fundamental question: \emph{even in highly simplified versions of UNO, does a unique Bayes-optimal action typically exist for a rational agent?}

Our approach is to formalize two-player UNO as a POMDP from the perspective of Player~1. Given an observation (own hand, opponent hand size, deck size, and public cards), we enumerate all latent world states that are consistent with that observation. Under a given value function we then compute the exact expected return of every legal action by averaging over all such worlds.

Experiments on reduced decks reveal that states with a strictly dominant action are uncommon: most belief states admit several actions with nearly identical expected value. In other words, even when actions are evaluated under a God’s-eye view of the underlying hidden state, UNO rarely exhibits a clearly optimal play. 



\section{POMDP Formulation of Two-Player UNO}

We formalize a modified two-player UNO environment as a partially observable
Markov decision process (POMDP). This section defines the latent state,
the player observation, the transition model, legal actions, and reward
structure.

\subsection{{State Space}}

Let the full card deck be
\[
D = H_1 \cup H_2 \cup D_g \cup P.
\]

The underlying environment state is
\[
S = (H_1, H_2, D_g, P, P_t, G_o),
\]
where
\begin{itemize}
    \item $H_1, H_2$: hands of Player~1 and Player~2,
    \item $D_g$: remaining draw deck (unordered set),
    \item $P$: previously played cards,
    \item $P_t = (\text{COLOR}, \text{VALUE})$: the top card,
    \item $G_o \in \{\text{Active}, \text{GameOver}\}$: terminal flag.
\end{itemize}

Wild and +4 cards are represented as
\[
P_t = (\text{CHOSEN\_COLOR},\ \text{VALUE}).
\]

\subsection{Action Space}

Player~1’s action space is
\[
A = \{X_1\} \cup \{Y_n : n \in \{1,2,4\}\}.
\]

\subsubsection{Play Action}

A play action is defined as
\[
X_1 = (\text{COLOR},\ \text{VALUE}),
\]
with special cases:
\[
(\text{BLACK},\ \text{WILD},\ \text{CHOSEN\_COLOR}),
\]
\[
(\text{BLACK},\ +4,\ \text{CHOSEN\_COLOR}).
\]

\subsubsection{Draw Action}

Drawing is represented as
\[
Y_n = \text{draw } n \text{ cards},
\]
where
\begin{itemize}
    \item $n = 1$ if no legal action exists,
    \item $n = 2$ or $n = 4$ if forced by +2 or +4 cards.
\end{itemize}


In this version of UNO, a player may not immediately play their turn after drawing any cards.

\subsection{Observation Model}

Player~1 observes
\[
O = (H_1,\ |H_2|,\ |D_g|,\ P,\ P_t,\ G_o),
\]
which includes its own hand, public card information, and deck size.
Opponent hand contents and the deck order remain hidden.

The observation function is deterministic:
\[
O(o \mid s', a) =
\begin{cases}
1, & o = (H_1', |H_2'|, |D_g'|, P', P_t', G_o'), \\
0, & \text{otherwise}.
\end{cases}
\]

Special events such as chosen colors for WILD/+4 and turn-skips induced
by SKIP/REVERSE are directly observable.

\subsection{Legal Moves}

A card $c$ is legally playable if
\[
\text{LEGAL}(P_t)
=
\left\{
c \in D :
\begin{array}{l}
c(\text{COLOR}) = P_t(\text{COLOR})\ \lor\\
c(\text{VALUE}) = P_t(\text{VALUE})\ \lor\\
c(\text{COLOR}) = \text{BLACK}
\end{array}
\right\}.
\]

\subsection{Transition Function}

State transitions follow:
\[
\resizebox{0.88\linewidth}{!}{$
T(s' \mid s, a) =
\begin{cases}
1,
&
\begin{aligned}
&a = X_1,\ X_1 \in H_1 \cap \text{LEGAL}(P_t), \\
&H_1' = H_1 \setminus \{X_1\}, \\
&P' = P \cup \{X_1\}, \\
&P_t' =
  \begin{cases}
  (\text{CHOSEN\_COLOR},\ \text{VALUE}), & \text{wild or +4},\\
  X_1, & \text{otherwise},
  \end{cases} \\
&G_o' =
  \begin{cases}
  \text{GameOver}, & |H_1'| = 0,\\
  \text{Active}, & \text{otherwise},
  \end{cases}
\end{aligned}
\\[1.2em]

\dfrac{1}{|D_g|},
&
\begin{aligned}
&a = Y_1,\ \text{no legal plays},\\
&H_1' = H_1 \cup \{c\},\ c \in D_g,\\
&D_g' = D_g \setminus \{c\};
\end{aligned}
\\[1.2em]

\dfrac{1}{\binom{|D_g|}{n}},
&
\begin{aligned}
&a = Y_n,\ P_t(\text{VALUE}) \in \{+2,+4\},\\
&H_1' = H_1 \cup J_n,\ |J_n| = n,\ J_n \subset D_g,\\
&D_g' = D_g \setminus J_n,\ P_t' = P_t;
\end{aligned}
\\[1.2em]

0,
&\text{otherwise or } G_o = \text{GameOver}.
\end{cases}
$}
\]


\subsection{Reward Function}

\[
R(s', s, a) =
\begin{cases}
+100, & |H_1'| = 0,\\
-100, & |H_2'| = 0,\\
0, & \text{otherwise}.
\end{cases}
\]

\subsection{Heuristic Value Function}

For reduced UNO experiments we additionally use the heuristic
\[
V(s) = |H_2| - |H_1|.
\]

\section{Particle Filtering for Belief State Approximation}

Since the state space is too large to enumerate all possible states, we use
particle filtering (a form of Bayesian filtering) to maintain an approximate
belief distribution. This approach generates particles dynamically at runtime,
caching them per game state for efficiency.

\subsection{Player 1's Perspective}

Player~1 observes:
\[
O = (H_1, |H_2|, |D_g|, P, P_t, G_o)
\]

Player~1 maintains beliefs about Player~2’s hidden hand $H_2$ using particles.
Each particle represents a possible full game state:
\[
s^{(i)} = (H_1, H_2^{(i)}, D_g^{(i)}, P, P_t, G_o).
\]

After Player~1 acts, environment dynamics propagate Player~2’s turn. Player~1
then observes the new public state and updates particles according to the
opponent’s behavior.

\subsection{Particle-Based Belief Representation}

Instead of maintaining a full belief distribution
\[
b(s) = \Pr(s \mid o, \text{history}),
\]
we use a particle set:
\[
\mathcal{P} = \{ (s^{(i)}, w^{(i)}) \}_{i=1}^{N},
\]
where:
\begin{itemize}
    \item $N$ is the number of particles (e.g., 1000–5000),
    \item $s^{(i)}$ is a sampled game state,
    \item $w^{(i)} \ge 0$ are importance weights,
    \item $\sum_i w^{(i)} = 1$.
\end{itemize}

Key insight: Only $H_2^{(i)}$ needs to be sampled, as $H_1$, $P$, $P_t$, and
$G_o$ are fully observed.

\subsection{Dynamic Particle Generation}

Given an observation
\[
O = (H_1, |H_2|, |D_g|, P, P_t, G_o),
\]
we compute the unknown card pool:
\[
L = D \setminus H_1 \setminus P.
\]

We enforce:
\[
|L| = |H_2| + |D_g|.
\]

For each particle $i = 1,\ldots,N$:
\begin{enumerate}
    \item Sample $H_2^{(i)}$ uniformly from all $\binom{|L|}{|H_2|}$ subsets.
    \item Set $D_g^{(i)} = L \setminus H_2^{(i)}$.
    \item Assign weight $w^{(i)} = 1/N$.
\end{enumerate}

If a user-provided deck size does not match $|L|-|H_2|$, it is automatically corrected.

\subsection{Bayesian Filtering Update}

When observation $o$ is received after action $a$, we apply:

\textit{1) Prediction step:}
\[
s'^{(i)} \sim T(\cdot \mid s^{(i)}, a).
\]

\textit{2) Observation update:}
\[
w'^{(i)} = w^{(i)} \cdot O(o \mid s'^{(i)}, a).
\]

\textit{3) Weight normalization:}
\[
w^{(i)} = \frac{w'^{(i)}}{\sum_j w'^{(j)}}.
\]

\textit{4) Resampling:  }

Define effective sample size:
\[
\text{ESS} = \frac{1}{\sum_i (w^{(i)})^2}.
\]
If $\text{ESS} < N/2$, resample $N$ particles and reset
$w^{(i)} = 1/N$.

\subsection{Specific Update Cases}

\subsubsection{Case 1: Opponent Plays a Card $z$}

Observation: opponent plays $z$, so $|H_2'| = |H_2| - 1$.

For each particle:
\begin{itemize}
    \item If $z \notin H_2^{(i)}$, set $w^{(i)} = 0$.
    \item If $z \in H_2^{(i)}$:
    \[
    H_2'^{(i)} = H_2^{(i)} \setminus \{z\},
    \]
    update $P$, $P_t$ accordingly, keep $D_g^{(i)}$ unchanged.
\end{itemize}

Normalize and resample afterward.

\subsubsection{Case 2: Opponent Draws a Card}

Observation: opponent has no legal moves and draws one card, so  
$|H_2'| = |H_2| + 1$.

For each particle:
\begin{itemize}
    \item If $H_2^{(i)}$ contains a legal move, set $w^{(i)} = 0$.
    \item Otherwise:
    \[
    c^{(i)} \sim \text{Uniform}(D_g^{(i)}),
    \]
    \[
    H_2'^{(i)} = H_2^{(i)} \cup \{c^{(i)}\},
    \]
    \[
    D_g'^{(i)} = D_g^{(i)} \setminus \{c^{(i)}\},
    \]
    \[
    w^{(i)} \leftarrow w^{(i)} \cdot \frac{1}{|D_g^{(i)}|}.
    \]
\end{itemize}

Normalize and resample.

\section{Experiment Setup}


Our project originally attempted to plan optimally in full two-player
UNO using POMDP methods. We first implemented two approximate planning
approaches:

\begin{itemize}
    \item \textbf{Particle Filter + Expectiminimax} on a simplified UNO
          deck containing only colored number cards.
    \item \textbf{Monte Carlo Tree Search (MCTS)} on the full UNO deck,
          using particle filtering to represent belief over hidden cards.
\end{itemize}

Both approaches yielded measurable but limited improvements over a
random baseline. Increasing particle counts or search depth produced
only marginal gains, suggesting that approximation quality alone did not
explain the difficulty of achieving consistently strong play.

\subsection{Reduced UNO Environment}

To enable exact analysis of belief states, we construct a drastically
reduced UNO deck in which all possible world states can be enumerated.
For example, the small deck contains only two colors and three numeric values:

\[
D_{\text{small}} =
\{(R,1), (R,2), (R,3),\ (G,1), (G,2), (G,3)\}.
\]

There are no action cards (SKIP, REVERSE, +2, +4, WILD), and card effects
are purely based on matching color or value. A typical scenario used in
our experiments is:

\[
H_1 = \{(R,1), (G,2)\},\quad |H_2| = 1,
\]
\[
P = \{(G,3)\},\quad P_t = (G,1),\quad |D_g| = 2.
\]

Given the player's observation
\[
O = (H_1,\ |H_2|,\ |D_g|,\ P,\ P_t,\ G_o),
\]
all hidden states \((H_2, D_g)\) consistent with the observation can be
enumerated exactly. For each legal action, its expected value is computed
by averaging over all such belief-consistent worlds.

This environment is small enough to evaluate the \emph{true}
Bayes-optimal Q-values, allowing us to determine whether a unique
optimal action exists in each scenario.

<ALEX AAND WESLY INSERT YOUR FINDINGS HERE>


\subsection{Full UNO with Monte Carlo Tree Search}





\section{Conclusion}

The conclusion goes here.

\bibliographystyle{IEEEtran}
\begin{thebibliography}{1}
\bibitem{IEEEhowto:kopka}
H.~Kopka and P.~W. Daly, \emph{A Guide to \LaTeX}, 3rd~ed.\hskip 1em plus
  0.5em minus 0.4em\relax Harlow, England: Addison-Wesley, 1999.
\end{thebibliography}

\end{document}



